\chapter{Learning as Infrastructure}

\chapterepigraph{%
  Learning is not filling a container.\\
  It is building a city —\\
  always half-finished, always expanding,\\
  always requiring maintenance.
}{Flyxion, \textit{Semantic Infrastructure}}

\sidenote{Connects to\\App.\,N on\\learning geometry;\\ Ch.\,23 on\\memory chains}

The accessibility framework gives us a precise way to understand learning:
it is not information transfer (pouring knowledge from teacher to student)
but \emph{infrastructure construction} (building the constraint structures
that make new knowledge accessible).

\section{Knowledge Accumulation}

Learning accumulates constraint modifications into a personal knowledge
infrastructure. Each new concept learned modifies the constraint set
in ways that make some future concepts easier to learn (more accessible)
and others harder (now inconsistent with established beliefs).

\begin{definition}{Knowledge Infrastructure}{knowledge-infra}
  An agent's \emph{knowledge infrastructure} at time $t$ is the pair
  $(\mathcal{C}(t), S_K(t))$: the current constraint set and the
  accessibility landscape it induces. The infrastructure determines
  what new concepts can be learned (those with high $S_K$ relative to
  the current constraint set) and at what cost.
\end{definition}

This explains why prerequisites matter. Learning calculus without algebra
is difficult not because calculus is intrinsically harder, but because
the accessibility of calculus concepts depends on the constraint
structures that algebra builds. Algebra makes calculus accessible;
without it, calculus concepts have near-zero accessibility in the
learner's knowledge infrastructure.

\section{Maintenance Costs}

Knowledge, like physical infrastructure, requires maintenance. Concepts
not visited decay (Chapter~23). Connections not reinforced weaken.
Beliefs not updated drift from the current state of knowledge.

The total maintenance cost of a knowledge infrastructure is:
\[
  \kappa_\text{maint}(\mathcal{C}, t) = \sum_{c \in \mathcal{C}}
  \text{decay}(c, t) \cdot \text{importance}(c)
\]
where $\text{decay}(c, t)$ is the rate at which constraint $c$
weakens without reinforcement. High-importance, high-decay constraints
(foundational concepts in fast-moving fields) require frequent maintenance.

\section{Forgetting as Accessibility Pruning}

Forgetting is often viewed negatively — as failure of the memory system.
The accessibility framework suggests a more nuanced view: forgetting is
\emph{accessibility pruning}, the removal of low-value constraints to
preserve the structural integrity of the knowledge infrastructure.

\begin{proposition}{Adaptive Forgetting}{adaptive-forgetting}
  In a knowledge infrastructure with finite maintenance capacity
  $B$ (total cognitive budget), optimal forgetting drops constraints
  $c$ with low $\text{value}(c) / \text{decay}(c)$ — those that
  are costly to maintain relative to their contribution to $S_K$.
\end{proposition}

This explains why we forget details (high decay, low value) while
retaining structure (low decay, high value). The structure of algebra
persists even when specific problem solutions are forgotten; the
structure enables reconstruction of details when needed.

\section{Reconstruction over Retention}

A mature knowledge infrastructure is not characterized by the retention
of large volumes of information but by the ability to reconstruct
needed information from structure. The expert mathematician who cannot
recall a specific theorem can re-derive it; the expert physician who
cannot recall a drug interaction can look it up and immediately integrate
it with their existing knowledge.

\begin{namedthm}{The Reconstruction Principle}
  A robust knowledge infrastructure prioritizes the retention of
  high-level constraint structures (schemas, frameworks, principles)
  over the retention of specific facts. Facts are reconstructible
  from structure; structure is not reconstructible from facts alone.
  The value of learning a fact is proportional to its contribution
  to structural constraint modifications.
\end{namedthm}

\begin{exercise}
  You are designing a curriculum for a 10-year, full-time mathematics
  education. Using the knowledge infrastructure model, what should
  be the primary goal of years 1--3 (foundation), 4--7 (expansion),
  and 8--10 (integration)? How does your design differ from a
  typical university mathematics curriculum?
\end{exercise}

\begin{exercise}
  Model the knowledge infrastructure of a first-year medical student
  versus a 20-year experienced physician. What are the key structural
  differences? Where does the physician's infrastructure support
  reconstruction that the student's does not?
\end{exercise}
